% from ega4-9-fr.tex
\oldpage[IV-3]{59}
Or, $f$ est de présentation finie (1.6.2, (v)), donc il résulte de (9.2.3) que
$Z$ est localement constructible (resp. ind-constructible) dans $Y$, donc
$Y-Z$ est localement constructible (resp. pro-constructible) dans $Y$. Comme
$g$ est de présentation finie, $g(Y-Z)$ est localement constructible
(resp. pro-constructible) dans $S$, en vertu du théorème de Chevalley (1.8.4)
(resp. de (1.9.5, (vii)))~; donc $E$ est localement constructible
(resp. ind-constructible) dans $S$.

\begin{remark}[9.2.5]
\label{IV.9.2.5-fr}
On notera que si $\mathbf P$ est une propriété de préschémas algébriques pour
laquelle la prop. (9.2.3) est vraie, alors $\mathbf P$ satisfait aussi à la
condition (9.2.1, (ii))~: cela résulte en effet du fait que dans un espace
irréductible noethérien, un ensemble constructible est rare ou contient un
ouvert non vide $(\mathbf 0_{\mathrm{III}},9.2.3)$.
\end{remark}

\begin{proposition}[9.2.6]
\label{IV.9.2.6-fr}
Désignons par $\mathbf P$ une des propriétés suivantes d'un $k$-préschéma
algébrique $X$~:
\begin{enumerate}[label=(\roman*)]
  \item $X$ est vide.
  \item $X$ est fini sur $k$.
  \item $X$ est radiciel sur $k$.
  \item $\dim(X)$ appartient à une partie donnée $\Phi$ de l'ensemble
  $\overline{\mathbf Z}=\mathbf Z\cup\{-\infty\}$.
\end{enumerate}
Alors $\mathbf P$ est constructible.
\end{proposition}

Il est clair que (i) et (ii) sont des cas particuliers de (iv), en prenant
respectivement pour $\Phi$ l'ensemble $\{-\infty\}$ et l'ensemble
$\{-\infty,0\}$. On n'a donc qu'à démontrer (iii) et (iv). Dans chacun de ces
deux cas la condition (i) de (9.2.1) est remplie en vertu de (2.7.1, (xv)) et
(4.1.4). D'autre part, dans le cas (iii), la propriété $\mathbf P$ vérifie la
conclusion de (9.2.3) en vertu de (1.8.7)~; il reste donc à voir qu'il en est de
même dans le cas (iv). Cela va résulter de la proposition plus précise
suivante~:

\begin{proposition}[9.2.6.1]
\label{IV.9.2.6.1-fr}
Si $f:X\to S$ est un morphisme de présentation finie, la fonction
$s\mapsto\dim(f^{-1}(s))$ est localement constructible.
\end{proposition}

La question est locale sur $S$, donc on peut supposer que
$S=\operatorname{Spec}(A)$ est affine et prouver que pour tout $n$, l'ensemble
$E$ des $s\in S$ tels que $\dim(X_s)=n$ est constructible. Le même
raisonnement que dans (9.2.3) ramène au cas où $A$ est noethérien et intègre,
et il suffit alors de prouver le

\begin{corollary}[9.2.6.2]
\label{IV.9.2.6.2-fr}
Soient $S$ un préschéma intègre noethérien de point générique $\eta$,
$f:X\to S$ un morphisme de type fini. Alors il existe un voisinage de $\eta$
dans $S$ tel que la fonction $s\mapsto\dim(X_s)$ soit constante dans ce
voisinage.
\end{corollary}

Les images par $f$ des composantes irréductibles (en nombre fini) de $X$ qui
ne rencontrent pas $X_\eta$, sont contenues dans des parties fermées de $S$ ne
contenant pas $\eta$ (puisque $S$ est intègre
$(\mathbf 0_{\mathrm I},2.1.5)$), donc (en remplaçant $S$ par un voisinage
ouvert de $\eta$) on peut se borner au cas où toutes les composantes
irréductibles $X_i$ de $X$ rencontrent $X_\eta$~; désignons encore par $X_i$
le sous-préschéma fermé réduit de $X$ ayant $X_i$ pour espace sous-jacent~;
comme $\dim(X_s)=\sup_i(\dim((X_i)_s))$ (4.1.1), on peut se borner au cas où
$X$ est irréductible. Il existe alors un recouvrement fini $(U_j)$ de $X$ par
des

